\label{sec:discussion}
\subsection{Critical evaluation of the design}
\label{sec:dis-critical}
The project connects photographed handwriting, numerical recognition, persistent writer profiles and a working Cartesian gantry. The results support this central function across eleven writers, with quality depending on the handwriting and embedded classification time remaining the principal performance constraint. This evaluation relates those outcomes to the decisions explained in Design.

\subsubsection{Interpretation of results}
\label{sec:dis-interpretation}
\paragraph{Useful reproduction is demonstrated, with a clear operating range.}
The 84.41\% read-back character score is close to the 85\% target, but its engineering meaning becomes clearer when the eleven writers are separated. Seven exceed 85\% and four exceed 93\%; those four include both personal contributors and a database writer. Thus the method is capable of substantially better than the pooled target on demonstrated styles. The lower scores on other writers show that the remaining improvement concerns consistency across styles, rather than an inability of the gantry to produce readable text at all. Figure~\ref{fig:res-physical} grounds this distinction in physical output.

The association between components per character and legibility is consistent with the limitations of the chosen representation. A glyph library stores observed character shapes and assembles them into new contexts. In more joined handwriting, the appearance of a character can depend strongly on its neighbours, so a cut glyph and a generic connection need not reproduce the same local shape. The observed rank association of 0.845 supports investigating this mechanism, but eleven writers do not isolate joining from glyph coverage, input quality or reader preference. The corresponding identification association is close to zero; it shows that a writer can remain identifiable even where individual characters are difficult for the reader.

The 86.41\% line-identification score and correct attribution of all fourteen page pools are useful evidence that writer-specific structure survives the full physical process. They answer a different question from stylistic correlation. The new physical geometric study gives seven-feature shape correlation 0.697 and nine-feature correlation 0.713, both below 0.85. It also identifies useful retained structure: slant tracks at 0.991, while ascender extent and character width track at approximately 0.817 and 0.798. Word spacing, at 0.486, is a weaker channel. These values motivate improving layout and context while preserving the successful orientation and glyph information; they do not establish that layout is the sole source of loss.

The writer model also participates in candidate selection, so it cannot be treated as an independent aesthetic judge of that selection. The geometric study uses no neural classifier, although it still depends on its descriptor definitions and the different original and reproduced texts. Its separately standardised coefficient primarily describes cross-writer ordering, and the lower 0.611 absolute shape variant shows that agreement of feature magnitudes is more demanding. Automated read-back is likewise an instrument-dependent measure of legibility. It provides a repeatable comparison, while a blinded human reading and style assessment would test the human-facing interpretation of the proposal more directly. Neither the physical overlay nor a human rating is inferred from model confidence.

\paragraph{The classifier results support the chosen architecture.}
The selected joint reader achieves 94.06\% pooled character accuracy, 0.94 percentage points below the 95\% target; the NumPy-trained scratch reader achieves 94.39\% on the same evaluation. The scratch reader's public and photographed subsets score 94.39\% and 94.42\%, while the joint reader scores 94.12\% and 93.33\%. These results show that the line-recognition architecture works on both distributor-cut IAM lines and the project's photographed material. Similar averages do not prove that segmentation is harmless: different text, writers and image conditions may compensate in the mean, and the study does not measure the segmentation contribution separately.

The 98.08\% writer-identification point estimate and correct pooling of all eleven writer groups are particularly strong outcomes for the multi-writer objective. The selected 0.17 stroke width makes this result relevant to a pen plotter that cannot reproduce each contributor's original ink thickness. The three errors also define the limits of the confidence output: one incorrect classification has an 83.8\% posterior, so confidence can guide review but cannot guarantee an identity. The evaluation is closed set, and its photographed split shares acquisition conditions with the training material; a new session with the same writers would be a stronger deployment check.

The selected library reader is retained as the intended demonstration configuration, while the alternative reader result is reported alongside it. This is a configuration decision rather than evidence that the library route necessarily gives the higher text score. On these data the scratch reader is slightly higher; for writer identification the library-trained model is higher. Keeping both comparisons visible gives a more useful assessment of the implementation than a single claim that one training route is uniformly superior.

\paragraph{The physical timing result is a practical success.}
All 25 explicitly completed positive-character ODROID records satisfy 4~s per character, including generation and drawing. The worst request still has 5.39~s of allowance remaining. This demonstrates usable physical throughput across the recorded 43--378-character range, rather than relying on the G-code generator's nominal feed calculation. That distinction is important because pen transitions, path length, GPIO scheduling and generation time all contribute to the request.

The shortest block has the smallest worst-case margin. Generation work is incurred before the drawing begins, while the allowance grows directly with character count. The measured worst short job includes 26.5~s of synthesis and approximately 3.248~s per character of drawing. Treating those two terms as fixed overhead plus a constant drawing rate would put a break-even length near 35 characters, but this is only a local projection. Search effort, repair and geometry change with the text, so the result motivates a direct short-request test rather than a claim that all shorter requests fail. Coordinate clamping in some recorded jobs also prevents the timing pass being used as a substitute for the separate fidelity result.

\paragraph{The classification-time gap is measurable and actionable.}
The two ODROID classification requests exceed their allowances by 31.76~s and 28.94~s. Approximately two thirds of the combined time is spent in the networks, with the remaining third in image processing. This makes network execution the largest observed optimisation opportunity, while also showing that segmentation alone cannot explain the total overrun. A two-request profile identifies the large contributors; it does not establish a constant cost law for arbitrary pages.

The development-host experiment supplies a concrete alternative: equivalent library networks batched by page are 2.76 times faster in the recorded comparison, with the reader text unchanged. It is reasonable to investigate that route on the ODROID. It is not yet reasonable to report the transferred factor as an embedded timing measurement, because the processor, libraries, thread scheduling and memory behaviour affect it. Even the projection with unchanged image processing leaves the two complete requests above their allowances. The required engineering work is therefore a coordinated reduction of network and image-stage cost, while retaining the output checks.

\paragraph{Mechanical accuracy is promising, with a precise distinction in what was measured.}
Every measured square and triangle feature is within 0.5~mm of nominal, and the square's mean absolute side error is 0.175~mm. These are useful results for the calibrated belt-and-rail mechanism. The triangle's unequal sloped sides additionally suggest an alignment effect that a side-length-only square check cannot expose. The apparent $0.706^\circ$ shear diagnostic offers a direction for adjustment, conditional on the assumed height and other geometric errors.

Neither equal side lengths nor small side-length errors fix the figure's absolute placement on the page. A translated figure can retain the same lengths, and repeated drawings can vary even when one specimen is accurate. The dimensional result should therefore be credited for what it demonstrates: a physically realised coordinate-to-motion path with sub-0.5~mm length differences on the measured features. A repeated-word overlay and a referenced bed-coordinate grid would extend that evidence to the proposal's positional and repetition objectives.

\paragraph{Relationship to published approaches.}
The sequence reader follows the CNN--BiLSTM--CTC architecture discussed by K{\i}z{\i}l{\i}rmak and Yan{\i}ko{\u g}lu~\cite[pp.~4--5]{kizilirmak2023}. The project's measured CER is 5.94\% for the selected reader and 5.61\% for the scratch reader, obtained directly from Equation~\ref{eq:res-char}. This is evidence of a functioning implementation in the same task family. It is not a controlled ranking against a published test result: the distributor's validation partition, training mixture, checkpoint selection and decoding protocol are not a matched published test protocol. In particular, the positional word-match score here is not directly comparable with published WER.

The reproduction strategy has a closer practical relationship to example-based approaches. Haines et al. demonstrate synthesis using a small collection of handwritten pages and explicit extracted forms~\cite[PDF p.~4, Sec.~4]{haines2016}. That supports the use of a small personal dataset as a design rationale without establishing an identical output quality. Graves's sequence generator uses ordered pen-coordinate observations~\cite[p.~18]{graves2013}; this project starts with offline photographs and does not observe those temporal paths. Modern image generators can model richer context, including One-DM's use of a single reference at generation time, but that reference count must be distinguished from the corpus and hardware used to train the generator~\cite[p.~9]{dai2024}. These comparisons explain the selected scope: an inspectable geometric system built from the available data, with context modelling as its principal reproduction trade-off.

\subsubsection{Critical evaluation}
\label{sec:dis-choices}
\paragraph{Separating the implementation contribution from checkpoint selection.}
The project implements convolution, normalisation, recurrent operations, CTC and optimiser updates numerically, as documented in Section~\ref{sec:des-implementations}. It also has a framework route for training the corresponding architecture. Choosing the library-trained writer for the demonstration does not remove the first-principles work; the numerical checks and trained-model comparison provide separate evidence for it. Conversely, the existence of a numerical trainer does not establish that every selected checkpoint came from that trainer.

Table~\ref{tab:dis-training} compares the saved models on common material. The numerical writer model exceeds the 95\% target as an observed line-level result, reaching 149/156 or 95.51\%, while the library-trained model reaches 153/156 or 98.08\%. On the 185-crop physical-output comparison the difference is larger, 80.54\% against 85.95\%. Both numerical and library checkpoints therefore learn useful writer distinctions, with the selected model providing the stronger measured writer result on these two populations. The 185-crop comparison is kept distinct from the 184 paired lines of QT~1.

\begin{table}[!htbp]\centering\small
\begin{tabular}{lrr}\toprule
\textbf{Common evaluation} & \textbf{Library trained} & \textbf{NumPy trained}\\\midrule
Writer: original validation lines & 153/156 (98.08\%) & 149/156 (95.51\%)\\
Writer: physical-output crops & 159/185 (85.95\%) & 149/185 (80.54\%)\\
Writer: original groups correct & 11/11 & 11/11\\
Writer: physical writer groups correct & 11/11 & 10/11\\
Reader: pooled character accuracy & 94.06\% & 94.39\%\\\bottomrule
\end{tabular}
\caption{Saved trained-model comparison. The reader row compares the joint and evaluated scratch checkpoints. The writer rows use matched crops and the same 0.17 preprocessing; these are comparisons of particular trained models, not an isolation of training precision or framework effects.}
\label{tab:dis-training}\end{table}

The supplied numerical-writer training log records 1\,031 training and 155 validation lines, a best displayed validation accuracy of 0.968 by epoch 41, and approximately 19--23~minutes per reported epoch. The reader excerpt records approximately 111--122~minutes per reported epoch and a resume from an existing numerical checkpoint. The later writer re-evaluation has 156 lines because the extraction and transcript correspondence used in that evaluation differs from the training-era population. Its 95.51\% result should therefore not be substituted for the 96.8\% training-validation figure or vice versa.

The reported float32 training choice is consistent with the need to limit the cost of these numerical runs, but this experiment does not isolate precision as the cause of an accuracy difference. Higher arithmetic precision does not, by itself, establish a higher trained-model classification score. A controlled precision study would require the same initial weights, sample order, preprocessing, optimiser settings and stopping rule. The existing fixed-parameter numerical checks support the implementation's arithmetic separately from these trained-model comparisons.

\paragraph{Line recognition addresses a real weakness in character isolation.}
The proposal described separating characters or connected groups before classification. The implemented reader instead recognises lines, while known-transcript CTC alignment recovers character boundaries during profile construction. This was a defensible response to connected handwriting: uncertain character cuts are not made a prerequisite for reading the text, and CTC supports training from complete transcripts~\cite[PDF pp.~2--3]{graves2006}. The cost is a larger sequence model and continued dependence on correct line ownership. The Results therefore qualify the final line-based design while retaining the functional contribution promised in the proposal.

\paragraph{The glyph library fits the data and exposes its own limitations.}
An explicit library permits size, slant, spacing and machine coordinates to be examined directly. It also allows the same stored writer to be used for different text without fitting a new generator for each request. These advantages are borne out by eleven functioning profiles and the different physical outputs. The price is a finite set of observed shapes, with limited information about how a writer changes a letter between neighbours. Multiple-character units or context-conditioned variants would address that loss more directly than increasing candidate count alone. They would also need reliable extraction and enough examples of the relevant combinations.

The first-principles and embedded-processing commitments of the proposal also constrain the interpretation of a library speed remedy. A framework benchmark is a useful engineering comparison; it does not by itself establish compliance with the original library restriction. A batched numerical implementation could retain that contribution, while compiled kernels for the identified hot loops offer a second route. Neither has an embedded compliance result in the present measurements. The practical workflow also separates off-device training from embedded operation. The supplied training records substantiate the numerical implementation and the cost of model preparation, while the proposal's intended embedded training mode is not demonstrated by those off-device runs. This distinction keeps the substantial implementation contribution and the final operating arrangement clear.

\paragraph{Economic feasibility.}
The author's current build-cost estimate is approximately R4\,000. The ODROID was borrowed at no purchase cost and will be returned after the project. This reduces the prototype's cash expenditure, but a replacement board would be part of a reproducible deployment cost. The ledger and subsequent purchases require reconciliation before a precise total is claimed; the later listed clipboard, paper, pen and estimated training electricity amount to R750 and may overlap the approximate estimate. The electricity amount is an author estimate, not an energy-meter result.

Using the available board allowed the project to prioritise the complete mechanical and image-processing system. The cost choice is relevant to the available compute, although the measurements do not show what a different board would have achieved. The 2.76-fold development-host comparison is stronger evidence for a software optimisation opportunity than an unsupported assertion that a more expensive platform would necessarily meet R4. A future cost comparison should include hardware, development effort and the measured quality retained by each acceleration route.

\subsubsection{Unsolved problems}
\label{sec:dis-unsolved}
The main reproduction problem is the variation between writers. Improvements should target the poorer cases without losing the high-performing printed and lightly joined styles. A useful next experiment would keep the reader, writer model and input text fixed while comparing current single-glyph paths with selected multi-character units. That would test the proposed context explanation directly. The new geometric study establishes a physical-output baseline; matching original and reproduced sentences and repeating capture sessions would better separate writer fidelity from content and acquisition effects.

The timing problem requires equivalent-output optimisation. The reduced-resolution experiment is valuable precisely because it measures the cost: it returns fewer lines on every comparable page and fails to segment two pages. It therefore cannot be accepted simply because its projected total may fall below an allowance. The better sequence is to preserve the reference output, measure the ODROID's hot loops and memory use, then compare a batched numerical route or selected compiled routines against that reference. The stopping criterion should be the complete request allowance together with unchanged transcripts and an agreed writer-output check.

The remaining mechanical work is a stronger positional qualification rather than an unsupported redesign. Frame squaring, belt consistency and pen-holder stiffness are plausible sources of geometric error, but the present length measurements do not rank them conclusively. Measuring repeated grid points, both square diagonals and the proposal's repeated word would identify where calibration is sufficient and where a mechanical adjustment is needed. A more controlled pulse scheduler and acceleration profile are sensible extensions of the motion design, but their benefit must be assessed on completed drawings rather than inferred from commanded feed.

The user has confirmed signed consent from the five contributors; institutional ethics status and the final material and expenditure records still need to be completed. These are documentation items separate from the measured engineering shortfalls. The model files, preprocessing width and stored writer mappings must also be carried into the demonstration as one identified configuration, so that a changed checkpoint does not silently change the meaning of the report's results.

\subsubsection{Strong points of the design}
\label{sec:dis-strengths}
The functioning chain is a principal strength: the system captures pages, recognises text, identifies writers, stores trajectories and produces real handwriting. Eleven styles are exercised in recognition and physical reproduction, seven exceed the reproduction character target, and all eleven writer-validation groups are identified correctly when pooled. These results substantiate the multi-writer objective.

An output can be followed through its source crop, normalised glyph, placement and emitted command. The independent G-code checks and physical measurements support this inspectable separation, allowing an improvement to be tested at a known interface while retaining the successful stages.

Operation-level checks and trained results support the numerical implementation: the numerical writer reaches 95.51\% on the common evaluation, and the numerical reader is competitive with the joint reader. The framework route supplies a practical reference and the stronger selected writer model. Together they support the implementation contribution and the final model choice.

\subsubsection{Expected failure conditions}
\label{sec:dis-failures}
Strong cursive connections, scarce character variants and text outside the available alphabet are the most relevant reproduction limits. A profile can preserve only the geometry it contains or the approximations selected for missing characters. Excessive layout extent can also cause clipping or coordinate clamping. A request should therefore be inspected in the preview and validated against the writing area before physical execution; a completed interpreter run alone is not proof of an unchanged trajectory.

For reading, unusually sparse pages, touching lines, page curvature and strong perspective changes can defeat the geometric preparation even when a cropped-line network is accurate. The Design's real capture examples illustrate these conditions. The writer classifier always chooses among its eleven represented identities, so an unfamiliar writer can receive a confident known-writer label. Short or ambiguous lines provide less evidence, while the present confidence sample does not establish a universally safe rejection threshold. Page pooling improves the observed group result but cannot recover information absent from every crop.

The gantry uses open-loop step counts. A missed step changes the physical position without an encoder measurement of that displacement. Limit switches and software abort handling provide useful boundary and stop functions, but depend on the installed switches, electrical circuit and running controller. The installed pen-motor circuit's lack of a flyback diode is a hardware limitation already identified in Section~\ref{sec:des-pen}; repeated reliable operation should not be treated as proof that the switching transient is harmless.

Finally, classification requests outside the tested sizes and very short writing requests have no demonstrated timing envelope. The measured results describe the recorded workload and platform, not all possible text, page sizes or concurrent users. This defines a practical laboratory prototype whose demonstrated successes can be repeated with an identified configuration, and whose extensions have specific measurements to target.

\subsection{Considerations in the design}
\label{sec:dis-considerations}
\subsubsection{Ergonomics}
The browser interface places text entry, captured-page review, writer selection and output preview on the operator's normal screen. Calibration and a writing-area indication reduce setup ambiguity, while the preview permits inspection before movement. Network operation is the selected arrangement; the ODROID also has a display interface~\cite[``External I/O'' section, unpaginated]{odroidn2ds}.

The page and pen still require manual positioning, and the motion area must remain accessible for loading while being kept clear during execution. A more complete user flow would confirm the paper restraint, valid calibration, selected profile and preview bounds before a job, and clearly distinguish generation, running, stopped and completed states. These are consequences of combining a software interface with a moving mechanism, rather than cosmetic interface additions.

\subsubsection{Health and safety}
The moving carriage, pulleys and belts present pinch points, and a protruding pen or tool can contact a hand in the working area. The installed end stops, limits, abort handling and controlled setup procedure reduce exposure during normal laboratory operation. The software stop is useful operational control; independent removal of motor power would provide protection if the controller or network became unresponsive. The driver and transistor temperature, wiring restraint and supply enclosure also require attention because the motor circuits carry the actual load currents.

The 24~V motor rail and 3.3~V logic are separate from the mains inputs of the purchased supplies. Low internal voltage does not establish mains insulation or enclosure quality. The supply labels and installed circuit should therefore be part of the final hardware record. A flyback path for the pen motor is a specific improvement supported by the switching analysis already given in Design~\cite[``Protection and Wiring Techniques'' section, unpaginated]{rohmlowside2026}; it is not described as a fitted protection component.

\subsubsection{Environmental impact}
The borrowed ODROID will be returned after the project. Reusable motors, rails, drivers and replaceable printed parts allow components to be repaired or used in other work. The printed material must be recorded before a particular recycling route can be specified.

Paper, pens and training electricity are the main recurring consumables identified by the author. Previewing jobs and testing geometry digitally can reduce wasted paper, while real physical tests remain necessary to establish performance. The reported R500 training-electricity allowance is a cost estimate; no total energy or environmental saving is calculated from it without power and runtime measurements. Electronics, plastics and remaining consumables should be separated and handled through appropriate local collection routes at end of life.

\subsubsection{Social and legal impact}
The intended benefit is personalised writing from digital text, including possible use by someone who cannot comfortably produce their own handwriting. The present system demonstrates the technical mechanism; accessibility or clinical benefit has not been evaluated with users. Reproducing another person's style also creates a risk of impersonation. Samples and profiles should be used with permission, and generated material should be identified as machine-produced where its origin matters.

The author confirms that all five personal contributors signed consent forms. Those signatures should be retained securely as supporting records; publishing them in the report would unnecessarily expose identifying material. Identifiable handwriting and writer profiles require consideration under South Africa's Protection of Personal Information Act, including minimal collection, a lawful processing basis, retention limits and security safeguards~\cite[PDF pp.~8--9, 11; ss.~10--11, 14 and 19]{popia2013}. Consent is one lawful basis and does not itself establish every aspect of compliance. The applicable database terms must also accompany use of the IAM material. The local-network server and stored profiles should have access restricted to the intended operator; the prototype is not presented as a publicly secured handwriting service.

\subsubsection{Ethics clearance}
Signed consent from the five contributors has been confirmed by the author. The institution's clearance or exemption decision is a separate matter and must be stated with its reference, where one exists.\designgap{Insert the university's confirmed ethics-clearance or exemption statement and reference. Participant consent is confirmed; do not infer institutional clearance from the signatures.}
\FloatBarrier
